% GENERATED FILE. Edit canonical lessons, then run npm run generate:books.
% canonical-source-hash: fnv1a64:87d12373ff769752
% canonical-lessons: GE-C119-schal, GE-C119-stiefel, GE-C119-rock, GE-C119-mantel, GE-C119-schuhe

\chapter{Scarf, Boots, Skirt, Coat, Shoes}
\label{ch:ge-scarf-boots-skirt-coat-shoes}
\hlchaptermodality{119}

\begin{chapteropening}
\textbf{By the end of this chapter:} \emph{I can say a scarf, boots, a skirt, a coat and shoes.}

Complete the last lesson of chapter 119: I can say a scarf, boots, a skirt, a coat and shoes.
\end{chapteropening}

\section[der Schal]{der Schal — a scarf}

\label{lesson:GE-C119-schal}



\begin{quote}
Before the new one: say the German for a tie, then the German for gloves.
\end{quote}

\subsection*{You'll want to know: der Schal}
\textbf{der Schal} — \textquotedblleft{}a scarf\textquotedblright{}.

\begin{cousinweb}
\textbf{Schal} $\leftarrow$ English \emph{shawl} $\leftarrow$ Persian \emph{šāl}, \textquotedblleft{}shawl, scarf\textquotedblright{} (EtymDB 2.1) — borrowed.
\end{cousinweb}

\begin{grammarlens}[title={What you've built}]
One more word for this chapter.
\end{grammarlens}

\subsection*{Your turn}
\begin{itemize}
\raggedright
  \item \emph{Say it:} \emph{der Schal}
  \item \emph{Say it:} \emph{der Schal}, once more
  \item \emph{Say it:} say \emph{der Schal}
  \item \emph{Recall:} say the German for a tie, then the German for gloves, then say \emph{der Schal} again
\end{itemize}

\begin{tcolorbox}[breakable,colback=teal!4,colframe=teal!35!black,title={Before you move on}]
What does der Schal mean? (\textquotedblleft{}A scarf\textquotedblright{}.) Say it once more.
\end{tcolorbox}
\section[die Stiefel]{die Stiefel — boots}

\label{lesson:GE-C119-stiefel}



\begin{quote}
Before the new one: say the German for gloves, then the German for a scarf.
\end{quote}

\subsection*{You'll want to know: die Stiefel}
\textbf{die Stiefel} — \textquotedblleft{}boots\textquotedblright{}.

\begin{cousinweb}
\textbf{Stiefel} $\leftarrow$ Old High German \emph{stival}, from Old French \emph{estivel} (EtymDB 2.1) — borrowed, probably from Latin \emph{aestīvāle}, a light summer shoe.
\end{cousinweb}

\begin{grammarlens}[title={What you've built}]
One more word for this chapter.
\end{grammarlens}

\subsection*{Your turn}
\begin{itemize}
\raggedright
  \item \emph{Say it:} \emph{die Stiefel}
  \item \emph{Say it:} \emph{die Stiefel}, once more
  \item \emph{Say it:} say \emph{die Stiefel}
  \item \emph{Recall:} say the German for gloves, then the German for a scarf, then say \emph{die Stiefel} again
\end{itemize}

\begin{tcolorbox}[breakable,colback=teal!4,colframe=teal!35!black,title={Before you move on}]
What does die Stiefel mean? (\textquotedblleft{}Boots\textquotedblright{}.) Say it once more.
\end{tcolorbox}
\section[der Rock]{der Rock — a skirt}

\label{lesson:GE-C119-rock}



\begin{quote}
Before the new one: say the German for a scarf, then the German for boots.
\end{quote}

\subsection*{You'll want to know: der Rock}
\textbf{der Rock} — \textquotedblleft{}a skirt\textquotedblright{}.

\begin{cousinweb}
\textbf{Rock} $\leftarrow$ Old High German \emph{roc}, \textquotedblleft{}tunic\textquotedblright{} $\leftarrow$ Proto-Germanic \emph{*hrukkaz}, \textquotedblleft{}robe, tunic\textquotedblright{} — inherited. (EtymDB 2.1) It first named a coat and later a skirt.
\end{cousinweb}

\begin{grammarlens}[title={What you've built}]
One more word for this chapter.
\end{grammarlens}

\subsection*{Your turn}
\begin{itemize}
\raggedright
  \item \emph{Say it:} \emph{der Rock}
  \item \emph{Say it:} \emph{der Rock}, once more
  \item \emph{Say it:} say \emph{der Rock}
  \item \emph{Recall:} say the German for a scarf, then the German for boots, then say \emph{der Rock} again
\end{itemize}

\begin{tcolorbox}[breakable,colback=teal!4,colframe=teal!35!black,title={Before you move on}]
What does der Rock mean? (\textquotedblleft{}A skirt\textquotedblright{}.) Say it once more.
\end{tcolorbox}
\section[der Mantel]{der Mantel — a coat}

\label{lesson:GE-C119-mantel}



\begin{quote}
Before the new one: say the German for boots, then the German for a skirt.
\end{quote}

\subsection*{You'll want to know: der Mantel}
\textbf{der Mantel} — \textquotedblleft{}a coat\textquotedblright{}.

\begin{cousinweb}
\textbf{Mantel} $\leftarrow$ Old High German \emph{mantal}, from Latin \emph{mantellum}, \textquotedblleft{}cloak\textquotedblright{} — borrowed, like English \emph{mantle}. EtymDB 2.1 has the Latin word.
\end{cousinweb}

\begin{grammarlens}[title={What you've built}]
One more word for this chapter.
\end{grammarlens}

\subsection*{Your turn}
\begin{itemize}
\raggedright
  \item \emph{Say it:} \emph{der Mantel}
  \item \emph{Say it:} \emph{der Mantel}, once more
  \item \emph{Say it:} say \emph{der Mantel}
  \item \emph{Recall:} say the German for boots, then the German for a skirt, then say \emph{der Mantel} again
\end{itemize}

\begin{tcolorbox}[breakable,colback=teal!4,colframe=teal!35!black,title={Before you move on}]
What does der Mantel mean? (\textquotedblleft{}A coat\textquotedblright{}.) Say it once more.
\end{tcolorbox}
\section[die Schuhe]{die Schuhe — shoes}

\label{lesson:GE-C119-schuhe}



\begin{quote}
Before the new one: say the German for a skirt, then the German for a coat.
\end{quote}

\subsection*{You'll want to know: die Schuhe}
\textbf{die Schuhe} — \textquotedblleft{}shoes\textquotedblright{}.

\begin{cousinweb}
\textbf{Schuhe} is the plural of \emph{Schuh} $\leftarrow$ Old High German \emph{scuoh} $\leftarrow$ Proto-Germanic \emph{*skōhaz}, the same word as English \emph{shoe}. (EtymDB 2.1)
\end{cousinweb}

\begin{grammarlens}[title={What you've built}]
One more word for this chapter.
\end{grammarlens}

\subsection*{Your turn}
\begin{itemize}
\raggedright
  \item \emph{Say it:} \emph{die Schuhe}
  \item \emph{Say it:} \emph{die Schuhe}, once more
  \item \emph{Say it:} say \emph{die Schuhe}
  \item \emph{Recall:} say the German for a skirt, then the German for a coat, then say \emph{die Schuhe} again
\end{itemize}

\begin{tcolorbox}[breakable,colback=teal!4,colframe=teal!35!black,title={Before you move on}]
What does die Schuhe mean? (\textquotedblleft{}Shoes\textquotedblright{}.) Say it once more.
\end{tcolorbox}
